Le filament de tungstène d'une ampoule à incandescence de \qty{\v{P}}{W} fonctionne à \qty{\v{T}}{K}. L'émissivité du tungstène à cette température est de \num{\v{em}}. On suppose que toute la puissance électrique est rayonnée par le filament et on néglige le rayonnement reçu de l'environnement.

\begin{enumerate}[label=\alph*)]
    \item Déterminez l'aire de la surface du filament.
    \item Déterminez la longueur d'onde du maximum d'émission. Pourquoi dit-on que ces ampoules sont surtout des « chaufferettes » ?
    \item Si le filament était un corps noir à la même température, quelle aire faudrait-il pour rayonner la même puissance ?
\end{enumerate}
